% !TEX program = lualatex
% Hoja de ejercicios: capítulo 3, teoría elemental de conjuntos

\documentclass[11pt,a4paper]{article}

\usepackage{fontspec}
\usepackage[spanish,es-nodecimaldot]{babel}
\usepackage[a4paper,margin=2.35cm,headheight=15pt]{geometry}
\usepackage{amsmath,amssymb,mathtools}
\usepackage{enumitem}
\usepackage{tabularx}
\usepackage{xcolor}
\usepackage{fancyhdr}
\usepackage{lastpage}
\usepackage{hyperref}

\setmainfont{Latin Modern Roman}
\setsansfont{Latin Modern Sans}
\setmonofont{Latin Modern Mono}

\definecolor{FdMBlue}{RGB}{20,55,100}
\definecolor{FdMLightBlue}{RGB}{235,241,248}
\definecolor{FdMGray}{RGB}{90,96,104}

\hypersetup{
  colorlinks=true,
  linkcolor=FdMBlue,
  urlcolor=FdMBlue,
  pdftitle={Ejercicios: teoría elemental de conjuntos},
  pdfauthor={Fundamentos de las Matemáticas I}
}

\pagestyle{fancy}
\fancyhf{}
\lhead{Fundamentos de las Matemáticas I}
\rhead{Teoría de conjuntos}
\cfoot{\thepage\ de \pageref{LastPage}}
\renewcommand{\headrulewidth}{0.4pt}

\setlist[enumerate]{leftmargin=2.4em,itemsep=0.35em,topsep=0.4em}
\setlist[itemize]{leftmargin=2em,itemsep=0.2em,topsep=0.3em}

\newcounter{ejercicio}
\newenvironment{ejercicio}[1][]{%
  \refstepcounter{ejercicio}%
  \par\medskip\noindent
  {\color{FdMBlue}\sffamily\bfseries Ejercicio \theejercicio%
  \if\relax\detokenize{#1}\relax\else\ --- #1\fi.}\quad
}{\par\medskip}

\newcommand{\respuesta}[1]{%
  \par\vspace{#1}\noindent
  \textcolor{FdMGray!45}{\rule{\linewidth}{0.3pt}}\par
}

\newcommand{\cajarecordatorio}[1]{%
  \begin{center}
    \colorbox{FdMLightBlue}{%
      \begin{minipage}{0.91\linewidth}
        \small #1
      \end{minipage}}
  \end{center}
}

\begin{document}

\begin{center}
  {\color{FdMBlue}\sffamily\bfseries\LARGE Hoja de ejercicios}\par
  \vspace{0.2em}
  {\sffamily\large Capítulo 3: teoría elemental de conjuntos}\par
  \vspace{0.7em}
  \rule{\linewidth}{0.6pt}
\end{center}

\noindent
\begin{tabularx}{\linewidth}{@{}lX@{\hspace{1.5em}}lX@{}}
  \textbf{Nombre y apellidos:} & \hrulefill &
  \textbf{Fecha:} & \hrulefill
\end{tabularx}

\vspace{0.8em}
\noindent
En los ejercicios de demostración deben declararse los conjuntos de referencia,
las variables y las hipótesis. Las igualdades entre conjuntos se justificarán
mediante doble inclusión o mediante equivalencias de pertenencia.

\cajarecordatorio{
  \textbf{Notación básica.}
  \(x\in A\) expresa pertenencia y \(A\subseteq B\) expresa inclusión.
  Si \(A,B\subseteq U\), se usarán \(A\cup B\), \(A\cap B\),
  \(A\setminus B\) y \(A^c=U\setminus A\). Además,
  \(A\times B\) denota el producto cartesiano y \(\mathcal P(A)\) el
  conjunto de todos los subconjuntos de \(A\).
}

\section*{A. Descripción de conjuntos y pertenencia}

\begin{ejercicio}[Extensión y comprensión]
Escribe por extensión, siempre que sea posible, los siguientes conjuntos y
determina su número de elementos:
\begin{enumerate}[label=\alph*)]
  \item \(A=\{n\in\mathbb N:n<8\}\).
  \item \(B=\{n\in\mathbb Z:-3\leq n<4\}\).
  \item \(C=\{n\in\mathbb N:n\text{ divide a }24\}\).
  \item \(D=\{x\in\mathbb R:x^2=4\}\).
  \item \(E=\{x\in\mathbb R:x^2+1=0\}\).
\end{enumerate}
Después describe \(\{0,3,6,9,12\}\) por comprensión de dos formas distintas.
\respuesta{5.0cm}
\end{ejercicio}

\newpage

\begin{ejercicio}[Colecciones bien definidas]
Decide cuáles de las siguientes descripciones determinan un conjunto sin
ambigüedad. En los casos problemáticos, explica qué información falta y propón
una reformulación precisa.
\begin{enumerate}[label=\alph*)]
  \item Los números enteros próximos a \(0\).
  \item Los números primos menores que \(100\).
  \item Los estudiantes altos de una clase.
  \item Los puntos del plano cuya distancia al origen es \(1\).
  \item Los libros interesantes de una biblioteca.
\end{enumerate}
\respuesta{4.4cm}
\end{ejercicio}

\newpage

\begin{ejercicio}[Pertenencia]
Sean
\[
  A=\{0,1,\{1\},\{0,1\},\varnothing\},
  \qquad B=\{\varnothing,\{\varnothing\}\}.
\]
Determina el valor de verdad de cada afirmación y justifica las que sean falsas:
\begin{enumerate}[label=\alph*)]
  \item \(1\in A\). \qquad
        \(\{1\}\in A\). \qquad
        \(\{1\}\subseteq A\).
  \item \(0\subseteq A\). \qquad
        \(\varnothing\in A\). \qquad
        \(\varnothing\subseteq A\).
  \item \(\varnothing\in B\). \qquad
        \(\{\varnothing\}\in B\). \qquad
        \(\{\varnothing\}\subseteq B\).
  \item \(B\in A\). \qquad \(B\subseteq A\).
\end{enumerate}
\respuesta{7.0cm}
\end{ejercicio}

\begin{ejercicio}[Igualdad por comprensión]
Decide cuáles de las siguientes parejas describen el mismo conjunto. Justifica
la respuesta sin limitarte a enumerar ejemplos.
\begin{enumerate}[label=\alph*)]
  \item \(\{2k:k\in\mathbb Z\}\) y \(\{n\in\mathbb Z:n\text{ es par}\}\).
  \item \(\{x\in\mathbb R:x^2=1\}\) y \(\{1\}\).
  \item \(\{n^2:n\in\mathbb Z\}\) y \(\{m\in\mathbb N:m\geq0\}\).
  \item \(\{3k+1:k\in\mathbb Z\}\) y
        \(\{n\in\mathbb Z:n-1\text{ es múltiplo de }3\}\).
\end{enumerate}
\respuesta{5.0cm}
\end{ejercicio}

\newpage
\section*{B. Inclusión, igualdad y conjunto vacío}

\begin{ejercicio}[Relaciones de inclusión]
Sean \(A=\{1,2\}\), \(B=\{1,2,3\}\), \(C=\{2,3\}\) y
\(D=\{1,2,\{1,2\}\}\). Determina todas las relaciones de inclusión que se
cumplen entre dos conjuntos distintos de la lista \(A,B,C,D\).
\respuesta{5.5cm}
\end{ejercicio}

\begin{ejercicio}[Formalización de la inclusión]
Escribe mediante cuantificadores y conectores lógicos:
\begin{enumerate}[label=\alph*)]
  \item \(A\subseteq B\).
  \item \(A\nsubseteq B\).
  \item \(A=B\).
  \item \(A\neq B\).
  \item \(A\subseteq B\subseteq C\).
\end{enumerate}
Explica qué tipo de testigo permite demostrar \(A\nsubseteq B\).
\respuesta{5.3cm}
\end{ejercicio}

\newpage

\begin{ejercicio}[Propiedades de la inclusión]
Demuestra a partir de la definición que, para cualesquiera conjuntos \(A,B,C\):
\begin{enumerate}[label=\alph*)]
  \item \(A\subseteq A\).
  \item Si \(A\subseteq B\) y \(B\subseteq C\), entonces \(A\subseteq C\).
  \item Si \(A\subseteq B\) y \(B\subseteq A\), entonces \(A=B\).
\end{enumerate}
\respuesta{8.0cm}
\end{ejercicio}

\begin{ejercicio}[El conjunto vacío]
\begin{enumerate}[label=\alph*)]
  \item Demuestra que \(\varnothing\subseteq A\) para todo conjunto \(A\).
  \item Explica por qué la prueba no afirma que exista un elemento de
        \(\varnothing\).
  \item Determina \(\mathcal P(\varnothing)\) y
        \(\mathcal P(\{\varnothing\})\).
  \item Decide si \(\varnothing\in\mathcal P(\varnothing)\) y si
        \(\varnothing\subseteq\mathcal P(\varnothing)\).
\end{enumerate}
\respuesta{6.0cm}
\end{ejercicio}

\newpage

\begin{ejercicio}[Cadenas de inclusiones]
Sean \(A,B,C\) conjuntos tales que \(A\subseteq B\subseteq C\). Decide cuáles
de las siguientes afirmaciones se deducen necesariamente. Demuestra las
verdaderas y da un contraejemplo para las falsas.
\begin{enumerate}[label=\alph*)]
  \item \(A\cap C=A\).
  \item \(A\cup C=C\).
  \item \(B\setminus A\subseteq C\setminus A\).
  \item \(C\setminus B\subseteq C\setminus A\).
  \item \(C\setminus A=(C\setminus B)\cup(B\setminus A)\).
\end{enumerate}
\respuesta{8.3cm}
\end{ejercicio}

\newpage
\section*{C. Operaciones con conjuntos}

\begin{ejercicio}[Cálculo con conjuntos finitos]
Sean
\[
 U=\{1,2,3,4,5,6,7,8,9,10\},\quad
 A=\{1,2,3,5,8\},\quad B=\{2,4,6,8,10\},\quad C=\{1,4,7,10\}.
\]
Calcula:
\begin{enumerate}[label=\alph*)]
  \item \(A\cup B\), \(A\cap B\), \(A\setminus B\) y \(B\setminus A\).
  \item \(A^c\), \(B^c\) y \(C^c\), respecto de \(U\).
  \item \(A\cap(B\cup C)\) y \((A\cap B)\cup(A\cap C)\).
  \item \((A\cup B)^c\) y \(A^c\cap B^c\).
  \item \((A\setminus B)\cup(B\setminus A)\).
\end{enumerate}
\respuesta{6.2cm}
\end{ejercicio}

\newpage

\begin{ejercicio}[El universo importa]
Sea \(A=\{1,2,3\}\). Calcula el complementario de \(A\) respecto de cada uno
de los universos siguientes:
\[
  U_1=\{1,2,3,4,5\},\qquad U_2=\mathbb N,
  \qquad U_3=\mathbb Z.
\]
Explica por qué la notación \(A^c\) es ambigua si no se conoce el universo.
\respuesta{5.0cm}
\end{ejercicio}

\begin{ejercicio}[Operaciones con intervalos]
Considera, como subconjuntos de \(\mathbb R\),
\[
  A=[-2,3),\qquad B=(1,5],\qquad C=(-\infty,0].
\]
Expresa mediante intervalos:
\begin{enumerate}[label=\alph*)]
  \item \(A\cap B\), \(A\cup B\) y \(A\setminus B\).
  \item \(A\cap C\), \(B\cap C\) y \((A\cup B)\cap C\).
  \item \(A^c\), \(B^c\) y \((A\cap B)^c\), tomando \(U=\mathbb R\).
\end{enumerate}
\respuesta{6.2cm}
\end{ejercicio}

\newpage

\begin{ejercicio}[Disjunción]
Para cada afirmación, decide si es verdadera. Si lo es, demuéstrala; si no,
da un contraejemplo.
\begin{enumerate}[label=\alph*)]
  \item Si \(A\cap B=\varnothing\), entonces \(A\subseteq B^c\).
  \item Si \(A\subseteq B\), entonces \(A\cap B^c=\varnothing\).
  \item Si \(A\cap B=\varnothing\) y \(A\cap C=\varnothing\), entonces
        \(B\cap C=\varnothing\).
  \item Si \(A\cup B=U\), entonces \(A\cap B=\varnothing\).
\end{enumerate}
\respuesta{7.0cm}
\end{ejercicio}

\begin{ejercicio}[Ecuaciones conjuntistas]
Sean \(A,B\subseteq U\). Describe todos los conjuntos \(X\subseteq U\) que
satisfacen cada ecuación, o demuestra que no existe ninguno:
\begin{enumerate}[label=\alph*)]
  \item \(X\cup A=A\).
  \item \(X\cap A=A\).
  \item \(X\setminus A=X\).
  \item \(X\cap A=\varnothing\) y \(X\cup A=U\).
\end{enumerate}
\respuesta{6.5cm}
\end{ejercicio}

\newpage
\section*{D. Lógica y lenguaje conjuntista}

\begin{ejercicio}[Traducir pertenencia a lógica]
Para un elemento arbitrario \(x\in U\), expresa cada pertenencia usando solo
las proposiciones \(x\in A\), \(x\in B\), \(x\in C\) y conectores lógicos:
\begin{enumerate}[label=\alph*)]
  \item \(x\in(A\cup B)^c\).
  \item \(x\in A\setminus(B\cap C)\).
  \item \(x\in(A\setminus B)\cup(B\setminus A)\).
  \item \(x\in(A\cup B)\setminus C\).
  \item \(x\in(A\cap B^c)\cup(A^c\cap B)\).
\end{enumerate}
\respuesta{5.5cm}
\end{ejercicio}

\begin{ejercicio}[Traducir lógica a conjuntos]
Sean \(P(x)\), \(Q(x)\) y \(R(x)\) las proposiciones \(x\in A\), \(x\in B\) y
\(x\in C\), respectivamente. Describe por operaciones conjuntistas el conjunto
de elementos \(x\in U\) que satisfacen:
\begin{enumerate}[label=\alph*)]
  \item \(P(x)\wedge\neg Q(x)\).
  \item \(\neg(P(x)\vee Q(x))\).
  \item \((P(x)\vee Q(x))\wedge\neg R(x)\).
  \item \((P(x)\Rightarrow Q(x))\).
  \item Exactamente una de las proposiciones \(P(x)\) y \(Q(x)\).
\end{enumerate}
\respuesta{6.0cm}
\end{ejercicio}

\newpage

\begin{ejercicio}[Diferencia simétrica]
Se define la diferencia simétrica por
\[
  A\mathbin{\triangle}B=(A\setminus B)\cup(B\setminus A).
\]
\begin{enumerate}[label=\alph*)]
  \item Describe con palabras sus elementos y represéntala en un diagrama.
  \item Demuestra que
        \(A\mathbin{\triangle}B=(A\cup B)\setminus(A\cap B)\).
  \item Calcula \(A\mathbin{\triangle}B\) para
        \(A=\{1,2,3,4\}\) y \(B=\{3,4,5\}\).
  \item ¿Cuándo se cumple \(A\mathbin{\triangle}B=\varnothing\)? Justifica.
\end{enumerate}
\respuesta{7.3cm}
\end{ejercicio}

\begin{ejercicio}[Diagnóstico de expresiones]
Corrige cada expresión que no esté bien formada o cuyo significado no coincida
con la frase indicada:
\begin{enumerate}[label=\alph*)]
  \item «\(x\) es subconjunto de \(A\)»: \(x\in A\).
  \item «\(A\) es un elemento de \(B\)»: \(A\subseteq B\).
  \item «\(x\) no pertenece ni a \(A\) ni a \(B\)»:
        \(x\notin A\cap B\).
  \item «Todo elemento de \(A\) pertenece a \(B\)»:
        \(\exists x\,(x\in A\Rightarrow x\in B)\).
  \item «Algún elemento de \(A\) no pertenece a \(B\)»: \(A\subseteq B\).
\end{enumerate}
\respuesta{5.0cm}
\end{ejercicio}

\newpage
\section*{E. Producto cartesiano}

\begin{ejercicio}[Enumerar productos cartesianos]
Sean \(A=\{1,2\}\), \(B=\{a,b,c\}\) y \(C=\varnothing\). Escribe por
extensión \(A\times B\), \(B\times A\), \(A\times A\), \(A\times C\) y
\(C\times B\). Compara sus cardinales.
\respuesta{6.0cm}
\end{ejercicio}

\begin{ejercicio}[Pertenencia a un producto]
Sean \(A,B\) conjuntos y \(a,b,c,d\) objetos. Justifica las equivalencias:
\begin{enumerate}[label=\alph*)]
  \item \((a,b)\in A\times B\Longleftrightarrow a\in A\text{ y }b\in B\).
  \item \((a,b)=(c,d)\Longleftrightarrow a=c\text{ y }b=d\).
\end{enumerate}
Después explica por qué \((a,b)\in A\times B\) no implica en general que
\((b,a)\in A\times B\).
\respuesta{5.5cm}
\end{ejercicio}

\newpage

\begin{ejercicio}[Cardinal del producto]
Sean \(A\) y \(B\) conjuntos finitos con \(|A|=m\) y \(|B|=n\).
\begin{enumerate}[label=\alph*)]
  \item Explica por qué \(|A\times B|=mn\).
  \item Si \(|A\times B|=12\), indica todas las posibilidades para
        \((|A|,|B|)\) cuando ambos conjuntos son no vacíos.
  \item Si \(|A\times A|=49\), determina \(|A|\).
  \item ¿Qué puede afirmarse si \(A\times B=\varnothing\)? Demuéstralo.
\end{enumerate}
\respuesta{6.7cm}
\end{ejercicio}

\begin{ejercicio}[Identidades con productos]
Demuestra o refuta, según corresponda:
\begin{enumerate}[label=\alph*)]
  \item \(A\times(B\cup C)=(A\times B)\cup(A\times C)\).
  \item \(A\times(B\cap C)=(A\times B)\cap(A\times C)\).
  \item \((A\cup B)\times C=(A\times C)\cup(B\times C)\).
  \item \((A\setminus B)\times C=(A\times C)\setminus(B\times C)\).
  \item \(A\times B=B\times A\).
\end{enumerate}
\respuesta{8.0cm}
\end{ejercicio}

\newpage
\section*{F. Conjunto de las partes}

\begin{ejercicio}[Enumerar subconjuntos]
Calcula explícitamente:
\begin{enumerate}[label=\alph*)]
  \item \(\mathcal P(\{a,b\})\).
  \item \(\mathcal P(\{1,2,3\})\).
  \item \(\mathcal P(\{\varnothing,1\})\).
\end{enumerate}
En cada caso comprueba que el número de elementos es una potencia de \(2\).
\respuesta{7.0cm}
\end{ejercicio}

\begin{ejercicio}[Pertenencia e inclusión en un conjunto de partes]
Sea \(A=\{1,2,3\}\). Decide si son verdaderas:
\begin{enumerate}[label=\alph*)]
  \item \(1\in\mathcal P(A)\). \qquad
        \(\{1\}\in\mathcal P(A)\).
  \item \(\{1\}\subseteq\mathcal P(A)\). \qquad
        \(\{\{1\}\}\subseteq\mathcal P(A)\).
  \item \(A\in\mathcal P(A)\). \qquad
        \(A\subseteq\mathcal P(A)\).
  \item \(\varnothing\in\mathcal P(A)\). \qquad
        \(\varnothing\subseteq\mathcal P(A)\).
\end{enumerate}
Justifica cuidadosamente cada respuesta.
\respuesta{6.5cm}
\end{ejercicio}

\newpage

\begin{ejercicio}[Contar subconjuntos]
Sea \(A\) un conjunto con \(8\) elementos. Determina:
\begin{enumerate}[label=\alph*)]
  \item El número total de subconjuntos de \(A\).
  \item El número de subconjuntos que contienen un elemento fijo \(a\in A\).
  \item El número de subconjuntos que no contienen \(a\).
  \item El número de subconjuntos que contienen dos elementos fijos distintos.
  \item El número de subconjuntos que contienen \(a\) pero no otro elemento
        fijo \(b\neq a\).
\end{enumerate}
Justifica cada recuento a partir de las decisiones de incluir o no cada elemento.
\respuesta{6.5cm}
\end{ejercicio}

\begin{ejercicio}[Inclusión y conjuntos de partes]
Demuestra que, para cualesquiera conjuntos \(A\) y \(B\),
\[
  A\subseteq B
  \quad\Longleftrightarrow\quad
  \mathcal P(A)\subseteq\mathcal P(B).
\]
Señala con claridad qué conjunto eliges en la implicación de derecha a izquierda.
\respuesta{7.5cm}
\end{ejercicio}

\newpage
\section*{G. Familias de conjuntos}

\begin{ejercicio}[Familia finita]
Sean
\[
 A_1=\{1,2,3,4\},\quad A_2=\{2,4,6\},\quad
 A_3=\{1,2,5,6\},\quad A_4=\{2,6,7\}.
\]
Calcula:
\[
 \bigcup_{i=1}^{4}A_i,\qquad
 \bigcap_{i=1}^{4}A_i,\qquad
 \bigcup_{i=2}^{3}A_i,\qquad
 A_1\cap\left(\bigcup_{i=2}^{4}A_i\right).
\]
\respuesta{5.2cm}
\end{ejercicio}

\begin{ejercicio}[Una familia aritmética]
Para cada \(n\in\mathbb N\setminus\{0\}\), define
\[
  A_n=\{k\in\mathbb Z:n\mid k\}.
\]
\begin{enumerate}[label=\alph*)]
  \item Describe \(A_1,A_2,A_3\) y \(A_6\).
  \item Calcula \(A_2\cap A_3\) y compáralo con \(A_6\).
  \item Calcula \(A_2\cup A_3\). ¿Coincide con algún \(A_n\)?
  \item Determina \(\bigcap_{n\geq1}A_n\).
\end{enumerate}
\respuesta{7.0cm}
\end{ejercicio}

\newpage

\begin{ejercicio}[Familias de intervalos]
Para \(n\geq1\), sean
\[
  I_n=\left(-\frac1n,\frac1n\right),
  \qquad J_n=[-n,n].
\]
Calcula y justifica:
\begin{enumerate}[label=\alph*)]
  \item \(\displaystyle\bigcap_{n=1}^{\infty}I_n\).
  \item \(\displaystyle\bigcup_{n=1}^{\infty}I_n\).
  \item \(\displaystyle\bigcap_{n=1}^{\infty}J_n\).
  \item \(\displaystyle\bigcup_{n=1}^{\infty}J_n\).
\end{enumerate}
\respuesta{7.0cm}
\end{ejercicio}

\newpage

\begin{ejercicio}[Cuantificadores en una familia]
Sea \((A_i)_{i\in I}\) una familia de conjuntos. Escribe con cuantificadores
las condiciones siguientes y niega correctamente cada una:
\begin{enumerate}[label=\alph*)]
  \item \(x\in\displaystyle\bigcup_{i\in I}A_i\).
  \item \(x\in\displaystyle\bigcap_{i\in I}A_i\).
  \item \(\displaystyle\bigcup_{i\in I}A_i\subseteq B\).
  \item \(B\subseteq\displaystyle\bigcap_{i\in I}A_i\).
\end{enumerate}
\respuesta{6.3cm}
\end{ejercicio}

\newpage

\begin{ejercicio}[De Morgan para familias]
Sea \((A_i)_{i\in I}\) una familia de subconjuntos de \(U\). Demuestra:
\[
  \left(\bigcup_{i\in I}A_i\right)^c
  =\bigcap_{i\in I}A_i^c,
  \qquad
  \left(\bigcap_{i\in I}A_i\right)^c
  =\bigcup_{i\in I}A_i^c.
\]
Haz explícito el cambio de «existe» a «para todo» al negar los cuantificadores.
\respuesta{8.0cm}
\end{ejercicio}

\section*{H. Identidades e inclusiones conjuntistas}

\begin{ejercicio}[Leyes básicas]
Demuestra mediante equivalencias de pertenencia:
\begin{enumerate}[label=\alph*)]
  \item \(A\cup\varnothing=A\) y \(A\cap U=A\).
  \item \(A\cup A=A\) y \(A\cap A=A\).
  \item \(A\cup A^c=U\) y \(A\cap A^c=\varnothing\).
  \item \((A^c)^c=A\).
\end{enumerate}
\respuesta{6.0cm}
\end{ejercicio}

\newpage

\begin{ejercicio}[Conmutatividad y asociatividad]
Demuestra:
\begin{enumerate}[label=\alph*)]
  \item \(A\cup B=B\cup A\) y \(A\cap B=B\cap A\).
  \item \((A\cup B)\cup C=A\cup(B\cup C)\).
  \item \((A\cap B)\cap C=A\cap(B\cap C)\).
\end{enumerate}
Explica por qué estas identidades permiten omitir paréntesis en uniones o
intersecciones sucesivas.
\respuesta{7.0cm}
\end{ejercicio}

\begin{ejercicio}[Distributividad]
Demuestra por doble inclusión las dos identidades distributivas:
\[
 A\cap(B\cup C)=(A\cap B)\cup(A\cap C),
\]
\[
 A\cup(B\cap C)=(A\cup B)\cap(A\cup C).
\]
\respuesta{8.0cm}
\end{ejercicio}

\newpage

\begin{ejercicio}[Leyes de De Morgan]
Sean \(A,B\subseteq U\). Demuestra:
\[
  (A\cup B)^c=A^c\cap B^c,
  \qquad
  (A\cap B)^c=A^c\cup B^c.
\]
Presenta una de las pruebas por doble inclusión y la otra mediante equivalencias
de pertenencia.
\respuesta{7.0cm}
\end{ejercicio}

\begin{ejercicio}[Identidades con la diferencia]
Demuestra las identidades:
\begin{enumerate}[label=\alph*)]
  \item \(A\setminus B=A\cap B^c\).
  \item \(A\setminus(B\cup C)=(A\setminus B)\cap(A\setminus C)\).
  \item \(A\setminus(B\cap C)=(A\setminus B)\cup(A\setminus C)\).
  \item \((A\setminus B)\setminus C=A\setminus(B\cup C)\).
\end{enumerate}
\respuesta{7.0cm}
\end{ejercicio}

\newpage

\begin{ejercicio}[Caracterizaciones de la inclusión]
Demuestra que son equivalentes las afirmaciones siguientes:
\[
  A\subseteq B,\qquad A\cap B=A,\qquad A\cup B=B,
  \qquad A\setminus B=\varnothing.
\]
Puedes probar una cadena de implicaciones o comparar cada igualdad con la
definición de inclusión.
\respuesta{8.0cm}
\end{ejercicio}

\begin{ejercicio}[Detectar una prueba incompleta]
Un estudiante intenta demostrar
\[
  A\cap(B\cup C)=(A\cap B)\cup(A\cap C)
\]
y escribe: «Sea \(x\in A\cap(B\cup C)\). Entonces \(x\in A\) y \(x\in B\) o
\(x\in C\). Por tanto, \(x\in(A\cap B)\cup(A\cap C)\), luego los conjuntos son
iguales».
\begin{enumerate}[label=\alph*)]
  \item Identifica la parte correcta y la conclusión que todavía no está
        justificada.
  \item Completa rigurosamente la demostración.
  \item Reescribe el paso lógico ambiguo usando paréntesis.
\end{enumerate}
\respuesta{6.5cm}
\end{ejercicio}

\newpage
\section*{I. Síntesis y problemas integradores}

\begin{ejercicio}[Verdadero o falso]
Decide si cada afirmación es verdadera para cualesquiera conjuntos. Demuestra
las verdaderas y refuta las falsas con un contraejemplo lo más sencillo posible.
\begin{enumerate}[label=\alph*)]
  \item \(A\setminus B\subseteq A\).
  \item \(A\subseteq A\cup B\).
  \item Si \(A\cup B=A\cup C\), entonces \(B=C\).
  \item Si \(A\cap B=A\cap C\), entonces \(B=C\).
  \item Si \(A\setminus B=B\setminus A\), entonces \(A=B\).
  \item \(\mathcal P(A\cap B)=\mathcal P(A)\cap\mathcal P(B)\).
  \item \(\mathcal P(A\cup B)=\mathcal P(A)\cup\mathcal P(B)\).
\end{enumerate}
\respuesta{9.0cm}
\end{ejercicio}

\newpage

\begin{ejercicio}[Reconstrucción de conjuntos]
En el universo \(U=\{1,2,3,4,5,6,7,8\}\), se sabe que
\[
 A\cap B=\{2,5\},\qquad A\setminus B=\{1,7\},\qquad
 B\setminus A=\{3,8\}.
\]
\begin{enumerate}[label=\alph*)]
  \item Determina \(A\), \(B\), \(A\cup B\) y \((A\cup B)^c\).
  \item Calcula \(A\mathbin{\triangle}B\).
  \item Determina todos los conjuntos \(C\subseteq U\) tales que
        \(A\cup B\subseteq C\) y \(|C|=7\).
\end{enumerate}
\respuesta{7.0cm}
\end{ejercicio}

\begin{ejercicio}[Partes y producto cartesiano]
Sean \(A=\{1,2\}\) y \(B=\{a,b\}\).
\begin{enumerate}[label=\alph*)]
  \item Enumera \(A\times B\) y calcula \(\mathcal P(A\times B)\).
  \item Compara \(\mathcal P(A\times B)\) con
        \(\mathcal P(A)\times\mathcal P(B)\): calcula sus cardinales y
        describe el tipo de sus elementos.
  \item Decide si ambos conjuntos pueden ser iguales y justifica la respuesta.
\end{enumerate}
\respuesta{7.5cm}
\end{ejercicio}

\newpage

\begin{ejercicio}[Una equivalencia completa]
Sean \(A,B\subseteq U\). Demuestra que son equivalentes:
\begin{enumerate}[label=\roman*)]
  \item \(A=B\).
  \item \(A\subseteq B\) y \(A^c\subseteq B^c\).
  \item \(A\mathbin{\triangle}B=\varnothing\).
  \item Para todo \(X\subseteq U\), se cumple
        \(A\cap X=B\cap X\).
\end{enumerate}
Organiza la prueba indicando explícitamente qué implicaciones demuestras.
\respuesta{9.0cm}
\end{ejercicio}

\newpage

\begin{ejercicio}[Problema final de demostración]
Sean \(A,B,C\subseteq U\). Demuestra la identidad
\[
  (A\cup B)\setminus C=(A\setminus C)\cup(B\setminus C)
\]
de dos maneras:
\begin{enumerate}[label=\alph*)]
  \item Mediante doble inclusión.
  \item Traduciendo la pertenencia a una equivalencia lógica.
\end{enumerate}
Después formula y demuestra la identidad dual obtenida al intercambiar unión e
intersección.
\respuesta{8.5cm}
\end{ejercicio}

\end{document}
